% Procedencia primaria (lectura de 18-09-2026):
% XROOT = output/EXTREMA_MEDIA_RAZON_NARRACION_Y_REVISION_20260916/ARTICULO/ES/source
% XROOT/libro/01_acoplamiento.tex:51-147,727-749.
% XROOT/sections/k_direccion_dimensional.tex:54-155.
% XROOT/nuclear/12.tex:9-88.
% VIROOT = output/AMPLIACION_FORMAL_LEAN_SERIE_20260916/delivery/USB_ES_EN_STAGE10_20260917/ES/PAQUETES/06_PARTICULAS_Y_MASAS_ES/documentacion_original/payload/source_es
% VIROOT/sections/05_registro_operador_masa.tex:11-119,193-294.
% VIROOT/sections/02_particula_persistencia.tex:210-235.
% XROOT/nuclear/07.tex:5-25,74-154: Hamiltoniano de transporte eliptico.
% XROOT/nuclear/11.tex:878-979: elevacion cotangente y accion discreta.
% Estatutos: cadena K->alpha->accion->masa = RESULTADO_RECUPERADO;
% carta 1+1+10 de eventos y criterio de entrelazamiento = FORMALIZACION_NUEVA;
% congruencia, espectro generalizado y control racional = pruebas completas aqui.

\section{Congruences positives, lecteurs d’action et dynamique hamiltonienne}
\label{sec:k-mass-bridge}

La relation entre le registre dodécaphasique et la Loi générale des masses possède une composition effective. APP produit les deux feuillets avec leurs résidus et quotients ; TRIT oriente la transition et détermine son régime ; TPK compose sélection, transport et actualisation en conservant route, retenue, frontière, survie et mémoire. Les représentations régionales et le registre à douze positions proviennent conjointement de l'état enrichi et de son prolongement compatible. La clôture numérique et la lecture matérielle partagent cette origine et conservent leurs opérations spécifiques.

L'origine de phase, l'orientation et la base mille étant déclarées, le registre est
\[
 K=(234,543,140,729,659,824,621,58,914,794,146,601).
\]
Le lecteur positionnel utilise
\begin{equation}
 \kappa_{\rm ag}(K)=\frac{\sum_{j=1}^{12}K_j1000^{12-j}}{1000^{12}-1},
 \qquad
 \alpha=\pi_{\rm HMT}+e_{\rm HMT}-\varphi_{\rm HMT}-4-
              \kappa_{\rm ag}(K).
 \label{eq:mass-k-alpha}
\end{equation}
Les valeurs du membre droit sont des représentations internes du même processus coinductif. La formule exprime le lecteur de clôture infinie compatible, avec la frontière de retenue qui le définit.

L'action reçoit ces représentations au moyen des lecteurs
\begin{align}
 H_5={}&\frac12\sqrt{1000\alpha/\varphi}-\alpha+
       \frac9{16}\alpha^2-\frac59\alpha^3+
       \frac7{48}\alpha^4-\frac1{54}\alpha^5,\notag\\
 D_A={}&\exp(-100\pi\alpha/9),\qquad
 \mathcal C_\pi=169\alpha^6+\frac{D_A\alpha^7}{1-D_A\alpha},\notag\\
 \eta_{\rm ret}={}&H_5-\frac{90}{\pi}\mathcal C_\pi,
 \qquad R_{\rm act}=\frac{H_5}{\eta_{\rm ret}}>0.
 \label{eq:mass-k-action}
\end{align}
Ici et dans les formules suivantes, \(\pi,\varphi,e\) abrègent les représentations HMT déjà construites. La norme déterminantale locale \(\Sigma_H=\varphi\alpha^{16}\) fixe la décade \(-34\), et la section \(\hbar_{\rm ret}=\eta_{\rm ret}10^{-34}\mathcal U_S\) appartient à la droite d'action déclarée. Le rapport \(R_{\rm act}\) conserve la relation entre les deux sections d'action lors du changement de leur base dimensionnelle commune.

La route électronique centrale fournit, par ses lecteurs de
déplacement et d’autocorrélation, le triplet \((80,54,6)\). Avec
\[
 \Delta_4=\frac{\pi-e\log\pi}{270}\left(1+\frac\pi{729}\right),
 \qquad
 b_e^{\rm alg}=\frac{\sqrt3}{4}
       \bigl[\exp(\varphi/\pi^2)-22\alpha^3\bigr](1+15\Delta_4),
\]
sa représentation matérielle contient
\begin{equation}
 \kappa_e=80-54\alpha+6\Delta_4,
 \qquad \varepsilon_e^\beta=b_e^{\rm alg}R_{\rm act}^{\kappa_e}.
 \label{eq:mass-k-electron}
\end{equation}
La composition \eqref{eq:mass-k-alpha}--\eqref{eq:mass-k-electron} identifie l'intervention concrète de \(K\) : sa lecture de clôture participe à \(\alpha\), et cette sortie agit ensuite dans le lecteur d'action et dans l'exposant de retour. Le choix central de l'électron relève de l'involution de la coordinatisation APP, dont l'unique point fixe est \((5,5)\). La centralité de cette cellule et son évaluation scalaire sont deux résultats reliés par la route et son registre.

\subsection{Coordonnées réversibles du registre d’événements matériels}

Le registre matériel assigne à une route \(\gamma\) douze événements
orientés \(\mathbf e(\gamma)=(e_0,\ldots,e_{11})\), un pour chaque
fenêtre nonadique du calendrier de cent huit positions. Le même
ordre temporel permet de réaliser ces événements dans le porteur à douze
coordonnées utilisé par \(K\). Cette identification requiert de conserver
l’origine et l’orientation du calendrier.

Soient \(\mathbf1\) le vecteur uniforme, \(P_{11}=I-\mathbf1\mathbf1^*/12\), et \(P_3\) le projecteur tridimensionnel de la coordinatisation incidentielle déclarée. La construction géométrique du registre détermine
\[
 u_K=P_3P_{11}K,\qquad
 \|u_K\|^2=\frac{6638585+2275584\sqrt5}{20}>0,
 \qquad u_K\perp\mathbf1,
\]
et le projecteur transversal
\[
 P_{10}(K)=P_{11}-\frac{u_Ku_K^*}{\|u_K\|^2}.
\]

\begin{proposition}[Décomposition matérielle relative à la direction de \(K\)]
Pour tout vecteur d'événements \(\mathbf e\), l'application
\begin{equation}
 \mathbf e\longmapsto
 \left(s_C=\mathbf1^*\mathbf e,\quad
       a_K=\frac{u_K^*\mathbf e}{\|u_K\|^2},\quad
       v_K=P_{10}(K)\mathbf e\right)
 \label{eq:mass-k-event-chart}
\end{equation}
est inversible sur son image et satisfait
\begin{equation}
 \mathbf e=\frac{s_C}{12}\mathbf1+a_Ku_K+v_K,
 \qquad
 \|\mathbf e\|^2=\frac{|s_C|^2}{12}
          +|a_K|^2\|u_K\|^2+\|v_K\|^2.
 \label{eq:mass-k-event-inverse}
\end{equation}
\end{proposition}
\begin{proof}
Les projecteurs \(\mathbf1\mathbf1^*/12\), \(u_Ku_K^*/\|u_K\|^2\) et \(P_{10}(K)\) sont deux à deux orthogonaux et leur somme est l'identité. Appliquer cette égalité à \(\mathbf e\) démontre l'inversion. L'orthogonalité des trois termes donne l'identité des normes.
\end{proof}

Cette coordinatisation conserve le vecteur complet des événements. Elle conserve donc également les fonctionnelles dépendant de leurs signes, de leurs sommes et de leurs corrélations, une fois effectuée la reconstruction \eqref{eq:mass-k-event-inverse}. Le registre des prédécesseurs des neuf, le décompte d'action et les coordonnées de tour demeurent dans la trajectoire enrichie : la coordinatisation précédente ne modifie que les coordonnées des douze événements. Le caractère est évalué sur la même signature récupérée. La construction fournit ainsi une représentation géométrique de l'information utilisée par la masse, avec une direction distinguée par \(K\) et dix composantes transversales.

La quantité d'information retenue par chaque représentation peut être démontrée exactement. Sur le domaine algébrique des registres, \(K'=K+\mathbf1\) vérifie
\[
 u_{K'}=u_K,\qquad P_{10}(K')=P_{10}(K),\qquad
 \kappa_{\rm ag}(K')-\kappa_{\rm ag}(K)=\frac1{999}.
\]
La dernière identité résulte de la sommation de la progression géométrique des douze poids positionnels. Si les représentations régionales sont maintenues fixes, la lecture de \(\alpha\) varie de \(-1/999\). Ce contrôle algébrique détermine la portée exacte de la direction : son projecteur conserve une polarisation, tandis que la clôture scalaire utilise en outre l'information uniforme et positionnelle du registre. L'admissibilité de \(K'\) comme autre historique terminal requiert sa propre construction APP--TRIT--TPK.

\subsection{Congruence positive et spectre de l'opérateur matériel}

Pour une multisection \(M\), soit \(\mathscr H_M\) l'espace fini muni d'une base orthonormale étiquetée par ses chemins. Le lecteur de retour \(r\in\{D,\Omega\}\) produit un opérateur autoadjoint \(B_M^r e_\gamma=\kappa_\gamma^r e_\gamma\). Le caractère raffiné, évalué sur les différences par rapport à l'électron, donne l'opérateur basal
\[
 M_M^{(0)}e_\gamma=
 m_e^{(0)}\chi_{\rm ref}(\sigma_\gamma-\sigma_e,
                  \eta_\gamma-\eta_e)e_\gamma.
\]
Les coefficients de tour et leur convergence relèvent du lecteur de chaque route. Leur représentation bêta est la congruence
\begin{equation}
 D_M^r=R_{\rm act}^{B_M^r/2},\qquad
 M_M^{\beta,r}=D_M^rM_M^{(0)}D_M^r.
 \label{eq:mass-k-congruence}
\end{equation}

\begin{proposition}[Invariants exacts de la congruence matérielle]
L’opérateur \(D_M^r\) est positif et inversible. La congruence
\eqref{eq:mass-k-congruence} conserve le rang et l’inertie ; en particulier,
elle conserve la positivité. Si \(M_M^{(0)}\) et \(B_M^r\) sont diagonaux
dans la base des routes, alors
\[
 m_\gamma^\beta=m_\gamma^{(0)}R_{\rm act}^{\kappa_\gamma^r}.
\]
Pour deux routes dans la même coordinatisation dimensionnelle, on obtient
\[
 \frac{m_Y^\beta}{m_X^\beta}=
 \chi_{\rm ref}(\sigma_Y-\sigma_X,\eta_Y-\eta_X)
 R_{\rm act}^{\kappa_Y^r-\kappa_X^r}.
\]
\end{proposition}
\begin{proof}
Le calcul fonctionnel donne \(D_M^r=\exp((\log R_{\rm act})B_M^r/2)\), de spectre strictement positif. Son inverse est l'exponentielle de l'opposé. Pour tout \(v\), la forme quadratique transformée vaut \((D_M^rv)^*M_M^{(0)}(D_M^rv)\). Le changement inversible de variable préserve les dimensions maximales des sous-espaces positifs, négatifs et nuls ; celles-ci déterminent l'inertie et le rang. Dans une base diagonale commune, les deux facteurs apportent conjointement \(R_{\rm act}^{\kappa_\gamma^r}\). Diviser deux valeurs annule la section dimensionnelle commune et utilise la loi multiplicative du caractère.
\end{proof}

La distinction entre transport d’une forme et modification d’un
opérateur relativement à une métrique fixe peut être vérifiée sans approximations.
Pour
\[
 M_0=\begin{pmatrix}1&0\\0&4\end{pmatrix},\qquad
 D=\begin{pmatrix}2&0\\0&1/2\end{pmatrix},\qquad
 M'=D^*M_0D=\begin{pmatrix}4&0\\0&1\end{pmatrix},
\]
le déterminant de \(D\) est un et les valeurs associées aux deux
routes s’échangent. Même le spectre non ordonné change si
\(M_0=I\) : il se transforme de \(\{1,1\}\) en \(\{4,1/4\}\).
Les deux contrôles préservent positivité, rang et inertie.

\begin{proposition}[Invariance spectrale avec métrique transportée]
Soient \(G_0>0\) une métrique et \(D\) inversible. Si
\[
 M'=D^*M_0D,\qquad G'=D^*G_0D,
\]
les paires \((M',G')\) et \((M_0,G_0)\) ont le même spectre
généralisé, avec les mêmes multiplicités.
\end{proposition}
\begin{proof}
On a
\[
 \det(M'-\lambda G')=|\det D|^2\det(M_0-\lambda G_0),
 \qquad (G')^{-1}M'=D^{-1}G_0^{-1}M_0D.
\]
La première égalité préserve les racines et leurs multiplicités ; la
seconde exhibe une similitude des opérateurs associés.
\end{proof}

Le flot diagonal positif induit par \(K\) possède donc deux réalisations mathématiques précises sur une forme matérielle. Lorsqu'il transporte simultanément la forme et sa métrique, il décrit le même objet dans une autre coordinatisation. Lorsqu'il agit sur la forme en conservant la métrique des routes, il produit une collection opératorielle pouvant modifier les masses relatives. Son identification physique est déterminée par le lecteur des routes et par l'action déclarée.

\subsection{Compatibilité du flux dodécaphasique avec le retour matériel}

Soit
\[
 A_K=(\log\rho_A)\operatorname{diag}(\widehat u_{K,1},\ldots,
       \widehat u_{K,12}),\quad
 \widehat u_K=u_K/\|u_K\|,\quad \rho_A=\pi/\varphi^2,
\]
de sorte que \(g_K(t)=e^{tA_K}\). La collection de retour matériel dans la multisection \(M\) a pour générateur
\[
 C_M^r=\tfrac12(\log R_{\rm act})B_M^r,\qquad
 S_M^r(t)=e^{tC_M^r}.
\]

\begin{proposition}[Transport de la congruence par un lecteur d'entrelacement]
Une application linéaire déclarée \(J_M:\mathbb C^{12}\to\mathscr H_M\) qui satisfait
\begin{equation}
 C_M^rJ_M=J_MA_K
 \label{eq:mass-k-intertwiner}
\end{equation}
satisfait
\[
 S_M^r(t)J_M=J_Mg_K(t),
\]
et, pour \(M_M(t)=S_M^r(t)M_M^{(0)}S_M^r(t)\),
\begin{equation}
 J_M^*M_M(t)J_M
 =g_K(t)^*\bigl(J_M^*M_M^{(0)}J_M\bigr)g_K(t).
 \label{eq:mass-k-pullback}
\end{equation}
\end{proposition}
\begin{proof}
L’identité \eqref{eq:mass-k-intertwiner} s’itère à chaque puissance
et se somme dans la série exponentielle convergente. Prendre les adjoints et
substituer dans la forme matérielle démontre
\eqref{eq:mass-k-pullback}.
\end{proof}

Le critère est fini et admet une épreuve de résistance directement opératorielle. Dans les bases diagonales déclarées, si \(a_j\) et \(c_\gamma\) sont les valeurs propres de \(A_K\) et \(C_M^r\), respectivement, on doit avoir
\[
 (c_\gamma-a_j)(J_M)_{\gamma j}=0
 \qquad\text{pour chaque }(\gamma,j).
\]
Chaque coefficient non nul du lecteur exige l'égalité correspondante des valeurs propres. La composition récupérée de \eqref{eq:mass-k-alpha}--\eqref{eq:mass-k-electron} et la coordinatisation réversible \eqref{eq:mass-k-event-chart} sont des résultats effectifs. L'identification supplémentaire de leurs deux flots requiert un lecteur \(J_M\) possédant la propriété \eqref{eq:mass-k-intertwiner} ; cette proposition détermine exactement ce que doit vérifier une réalisation qui la soutient.

\subsection{Transport elliptique et transformation régulière de Legendre}

La réalisation variationnelle du transport possède un antécédent explicite dans l'ellipse d'action. La paire angulaire déjà exprimée détermine \(\eta=\operatorname{artanh}(C^*/A)\), dans la chambre \(A\ne0\), \(|C^*/A|<1\). Les coordonnées \((Q,P)\), toutes deux de dimension racine d'action, réalisent la forme
\[
 \mathcal I_\eta(Q,P)=\frac12\bigl(e^{-\eta}Q^2+e^\eta P^2\bigr),
 \qquad \lambda=P\,dQ,\quad \Omega=dQ\wedge dP.
\]
Le transport entre formes, avec \(M_\eta=\operatorname{diag}(e^{\eta/2},e^{-\eta/2})\) et \(J_2=\left(\begin{smallmatrix}0&-1\\1&0\end{smallmatrix}\right)\), est \(M_{\eta'}e^{-\vartheta J_2}M_\eta^{-1}\). Une interpolation \(\eta(t)\), \(\vartheta(t)\), avec \(\dot\vartheta=\omega(t)>0\), possède l'hamiltonien récupéré de l'article X :
\begin{equation}
 H_{\rm tr}(Q,P,t)=\omega(t)\mathcal I_{\eta(t)}(Q,P)
                 +\frac{\dot\eta(t)}2QP.
 \label{eq:mass-k-hamiltonian}
\end{equation}
Le second terme provient de \(\dot M_\eta M_\eta^{-1}\). Les équations
\[
 \dot Q=\omega e^\eta P+\frac{\dot\eta}2Q,
 \qquad
 \dot P=-\omega e^{-\eta}Q-\frac{\dot\eta}2P
\]
conservent \(\mathcal I_{\eta(t)}(Q(t),P(t))\) : la variation explicite de \(\eta\) et la contribution du terme de connexion s'annulent. Ce résultat et l'identité \(P\dot Q-H_{\rm tr}=p\dot q-\omega(q^2+p^2)/2\), pour \(Q=e^{\eta/2}q\), \(P=e^{-\eta/2}p\), sont démontrés dans la section de conservation variationnelle de l'ouvrage sur l'extrême et moyenne raison holographique.

\begin{proposition}[Lagrangien régulier du transport elliptique]
Supposons \(\eta\in C^2\), \(\omega\in C^1\) et \(\omega>0\).
La transformation de Legendre par rapport à \(P\) est inversible et
produit
\begin{equation}
 L_{\rm tr}(Q,\dot Q,t)=
 \frac{\bigl(\dot Q-\dot\eta Q/2\bigr)^2}{2\omega e^\eta}
 -\frac{\omega e^{-\eta}}2Q^2.
 \label{eq:mass-k-lagrangian}
\end{equation}
\end{proposition}
\begin{proof}
Écrivons \(a=\dot\eta/2\), \(b=\omega e^\eta>0\) et
\(c=\omega e^{-\eta}\). Alors
\(H_{\rm tr}=bP^2/2+aQP+cQ^2/2\), et
\[
 \dot Q=\partial_PH_{\rm tr}=bP+aQ,
 \qquad P=\frac{\dot Q-aQ}{b}.
\]
En substituant cette expression dans \(P\dot Q-H_{\rm tr}\), les
termes contenant \(P\) se regroupent en
\(P(\dot Q-aQ)-bP^2/2=(\dot Q-aQ)^2/(2b)\), ce qui
démontre \eqref{eq:mass-k-lagrangian}. De plus,
\[
 \partial_{\dot Q}L_{\rm tr}=P,
 \qquad
 \partial_P^2H_{\rm tr}=\omega e^\eta>0,
 \qquad
 \partial_{\dot Q}^2L_{\rm tr}=\frac1{\omega e^\eta}>0.
\]
La transformation est régulière et ses équations d’Euler--Lagrange
équivalent aux équations hamiltoniennes précédentes.
\end{proof}

Cette élimination algébrique réunit le Hamiltonien déjà construit et
sa forme lagrangienne du second ordre. La régularité provient du
terme quadratique de la réalisation elliptique. Pour un relèvement
cotangent purement linéaire \(H_K(q,p)=p^{\mathsf T}A_Kq\), l’équation
\(\dot q=A_Kq\) est indépendante de \(p\) et le Hessien
\(\partial_p^2H_K\) est nul. Sa formulation variationnelle utilise
l’action du premier ordre \(\int p^{\mathsf T}(\dot q-A_Kq)\,dt\).
Le corpus contient également son antécédent discret
\(L_n=p_{n+1}^{\mathsf T}(q_{n+1}-U_nq_n)\), dont les équations
produisent le transport \(q_{n+1}=U_nq_n\) et le transport dual
\(p_{n+1}=U_n^{-\mathsf T}p_n\). Les deux formulations possèdent
des domaines et des opérateurs déterminés. L’interpolation temporelle et la
section d’action sont conservées comme parties de la réalisation physique
déclarée.
